%! Introducing Statistics: What Can We Learn from Data? — Unit 1, Lesson 1.0 — Guided Notes (Answer Key)
\documentclass[10pt]{article}
\usepackage{apstats-article}
\usepackage{apstats-key}   % pulls in -boxes; do NOT also load -boxes

% Vocab row — the vocabpar fix. \par on BOTH ends so a term label is never dragged
% onto the previous answer's dotted line. Same arity and same geometry as the blank's
% \vocabterm: two answer lines here where the blank reserves two write-lines, so the
% call sites below are byte-identical with the blank and cannot drift in height.
\newcommand{\vocabterm}[3]{%
  \par\noindent\textbf{\textcolor{navy}{#1:}}\\[1pt]%
  \ansline{#2}\\[1pt]\ansline{#3}\\[3pt]\par}

\begin{document}

\pageheader{Unit 1, Lesson 1.0}{Guided Notes --- Answer Key}

\vspace{0.15in}

\boxguard[20]
\begin{vocabbox}
\small
Fill in each term as we build it together.
\par\vspace{2pt}
\vocabterm{Statistics}
  {The science of collecting, organizing, analyzing, and drawing conclusions}
  {from data --- the whole process, not just the arithmetic.}
\vocabterm{Data}
  {Recorded values together with the context that gives them meaning.}
  {Values with no context are just numbers, not data.}
\vocabterm{Individual}
  {The person, animal, or object described by one row of a data set.}
  {Ask: ``who or what is one of these?''}
\vocabterm{Variable}
  {A characteristic recorded for each individual, which takes different}
  {values for different individuals. It is a column heading, not a value.}
\end{vocabbox}

\vspace{0.14in}

\boxguard
\begin{notesbox}{1. A Number Alone Tells You Nothing}
\small
The number \textbf{68} could be a test score, a temperature, a heart rate, or a height.
Until we say what it measures, it carries \ans{no} information.

\vspace{6pt}
\textbf{Essential Knowledge (VAR-1.A.1):} Numbers may convey meaningful information
\emph{when placed in} \ans{context}.

\vspace{8pt}
Before a number becomes data, we answer the \textbf{W questions}:

\vspace{6pt}
\renewcommand{\arraystretch}{1.7}
\begin{tabularx}{\linewidth}{@{} >{\bfseries\color{navy}}l X @{}}
\hline
Who?   & \ans{Seven students in this class.} \\
What?  & \ans{Their height.} \\
Units? & \ans{Inches.} \\
When \& Where? & \ans{Today, in this classroom.} \\
\hline
\end{tabularx}

\vspace{6pt}
Now reread the warm-up numbers as \emph{heights, in inches, of seven students in this
class.} The same seven numbers now answer a question.
\end{notesbox}

\vspace{0.14in}

\boxguard[30]
\begin{notesbox}{2. Individuals and Variables}
\small
A shelter records these facts about the dogs currently available:

\vspace{6pt}
\renewcommand{\arraystretch}{1.35}
\begin{tabularx}{\linewidth}{@{} l Y Y Y Y @{}}
\rowcolor{sky}
\textbf{Name} & \textbf{Breed} & \textbf{Age (yr)} & \textbf{Weight (lb)} & \textbf{Adopted?} \\
\hline
Biscuit & Beagle    & 3 & 24 & Yes \\
Momo    & Terrier   & 1 & 15 & No  \\
Ranger  & Shepherd  & 5 & 61 & Yes \\
Pepper  & Retriever & 2 & 55 & No  \\
\hline
\end{tabularx}

\vspace{10pt}
Each \textbf{row} is one \ans{individual}. \qquad
Each \textbf{column} is one \ans{variable}.

\vspace{8pt}
How many individuals are shown? \ans{4} \qquad
How many variables are recorded? \ans{5}

\vspace{8pt}
\textbf{Careful --- a variable is not a value.}\\[2pt]
``Beagle'' is a \ans{value} of the variable \ans{Breed}.

\vspace{8pt}
\textbf{Why does any of this matter?} If every dog weighed exactly the same, we would
need only one dog to know everything about weight. Because the values
\ans{differ}, we need data. That fact has a name:

\vspace{4pt}
\hspace*{1.2em}\textbf{\textcolor{navy}{Variability:}}\\[1pt]
\ansline{Values of a variable differ from one individual to the next.}
\end{notesbox}

\vspace{0.14in}

\boxguard[24]
\begin{notesbox}{3. What Makes a Question Statistical?}
\small
A \textbf{statistical question} passes \emph{both} checks:

\vspace{6pt}
\hspace*{1.2em}\textbf{Check 1.} Will the answers \ans{differ} across individuals?\\[4pt]
\hspace*{1.2em}\textbf{Check 2.} Would I have to \ans{collect data} to answer it?

\vspace{10pt}
Mark each question \textbf{S} (statistical) or \textbf{N} (not statistical).

\vspace{4pt}
\renewcommand{\arraystretch}{1.8}
\begin{tabularx}{\linewidth}{@{} X C{2.4cm} @{}}
\hline
\textbf{Question} & \textbf{S or N} \\
\hline
How many minutes did I sleep last night?          & \ans{N} \\
How many minutes do students at our school sleep? & \ans{S} \\
What is the weight of Ranger, the shepherd?       & \ans{N} \\
How do the weights of shelter dogs vary?          & \ans{S} \\
\hline
\end{tabularx}

\vspace{8pt}
Choose one question you marked \textbf{N}. What single change would make it
statistical?\\[4pt]
\ansline{``What is the weight of Ranger?'' asks about one dog, so it has one answer.}\\[1pt]\ansline{Change it to ``How do the weights of the shelter's dogs vary?'' --- now they differ.}
\end{notesbox}

\vspace{0.14in}

\boxguard[15]
\begin{practicebox}
\small
A gym records the resting heart rate, in beats per minute, of each of its 120 members
on the morning they join.

\vspace{8pt}
\textbf{(a)} The individuals are \ans{the 120 gym members}.

\vspace{6pt}
\textbf{(b)} The variable is \ans{resting heart rate}, measured in \ans{beats per minute}.

\vspace{6pt}
\textbf{(c)} Write one statistical question this data set could answer.\\[4pt]
\ansline{``How do the resting heart rates of the gym's members vary?'' The answers}\\[1pt]\ansline{differ from member to member, and data must be collected to find them.}

\vspace{6pt}
\textbf{(d)} Write one question about these members that is \emph{not} statistical.\\[4pt]
\ansline{``What is member \#7's resting heart rate?'' --- one member, one answer.}
\end{practicebox}

\end{document}
